\documentclass[11pt]{article}
\usepackage[margin=1in]{geometry}
\usepackage{amsmath,amssymb,amsthm}
\usepackage[hidelinks]{hyperref}

\newtheorem{theorem}{Theorem}[section]
\newtheorem{lemma}[theorem]{Lemma}
\newtheorem{corollary}[theorem]{Corollary}
\newtheorem{proposition}[theorem]{Proposition}
\theoremstyle{definition}
\newtheorem{definition}[theorem]{Definition}
\newtheorem{remark}[theorem]{Remark}

\newcommand{\C}{\mathbb{C}}
\newcommand{\R}{\mathbb{R}}
\newcommand{\Z}{\mathbb{Z}}
\newcommand{\N}{\mathbb{N}}
\newcommand{\Lam}{\Lambda}
\newcommand{\re}{\operatorname{Re}}
\newcommand{\im}{\operatorname{Im}}
\newcommand{\res}{\operatorname{Res}}
\newcommand{\lcm}{\operatorname{lcm}}
\newcommand{\G}{\mathcal{G}}
\newcommand{\RH}{\mathsf{RH}}
\newcommand{\GR}{\mathsf{G}}
\newcommand{\dd}{\,\mathrm{d}}

\title{The dyadic Goldbach rate and the Riemann hypothesis are one proposition\\
\large A self-contained proof of the biconditional}
\author{}
\date{}

\begin{document}
\maketitle

\begin{abstract}
\noindent
Let $\Lam$ be the von Mangoldt function, $R(N)=\sum_{a+b=N}\Lam(a)\Lam(b)$ the quantitative Goldbach count, and
$\G(t)=t^{2}\sum_{N}R(N)e^{-Nt}$ its normalised Laplace transform.
Write $\GR$ for the proposition
``for every $\varepsilon>0$, $\G(2t)-\G(t)=O_\varepsilon(t^{1/2-\varepsilon})$ as $t\downarrow 0$''
and $\RH$ for the Riemann hypothesis.
We prove $\GR\Leftrightarrow\RH$.
Hence any proof of either proposition is a proof of the other, and any refutation of either is a refutation of the other: the two have identical truth content.
Everything used is proved in this file from the base stated in Section~\ref{sec:base}: the real and complex numbers with the exponential function, Riemann integration of continuous functions, and the Lebesgue convergence and Fubini theorems.
On that base the file proves the needed complex analysis (Goursat, Cauchy, residues, the identity theorem, Schwarz, logarithms), Fourier analysis (Gaussian transform, inversion, Poisson summation), unique factorisation, the theory of $\Gamma$ (reflection formula, decay, Cahen--Mellin integral), and of $\zeta$ (Euler product, continuation, non-vanishing on $\re s=1$, functional equation), before the two directions of the theorem.
\end{abstract}

\section{Statement}\label{sec:statement}

\begin{definition}[Arithmetic objects]
For $n\ge 1$ let $\Lam(n)=\log p$ if $n=p^{k}$ for a prime $p$ and an integer $k\ge 1$, and $\Lam(n)=0$ otherwise.
Let $\psi(x)=\sum_{n\le x}\Lam(n)$ for real $x\ge 1$.
For $N\ge 2$ let
\[
R(N)=\sum_{\substack{a+b=N\\ a,b\ge 1}}\Lam(a)\Lam(b),
\]
the weighted ordered count of representations of $N$ as a sum of two prime powers.
For $t>0$ let
\[
A(t)=\sum_{n\ge 1}\Lam(n)e^{-nt},\qquad
G_R(t)=\sum_{N\ge 2}R(N)e^{-Nt},\qquad
\G(t)=t^{2}G_R(t),\qquad
r(t)=\G(2t)-\G(t).
\]
\end{definition}

All series converge absolutely for $t>0$ since $0\le\Lam(n)\le\log n\le n$ and $R(N)\le N^{2}\cdot N=N^{3}$.

\begin{lemma}\label{lem:square}
For every $t>0$, $G_R(t)=A(t)^{2}$; hence $\G(t)=(tA(t))^{2}$.
\end{lemma}
\begin{proof}
All terms are nonnegative, so by Tonelli
\[
A(t)^{2}=\sum_{a\ge1}\sum_{b\ge1}\Lam(a)\Lam(b)e^{-(a+b)t}=\sum_{N\ge2}e^{-Nt}\sum_{a+b=N}\Lam(a)\Lam(b).\qedhere
\]
\end{proof}

\begin{definition}[The two propositions]
\begin{itemize}
\item[$\GR$:] For every $\varepsilon\in(0,\tfrac12)$ there exist $C_\varepsilon>0$ and $t_\varepsilon>0$ such that
$|r(t)|\le C_\varepsilon\,t^{1/2-\varepsilon}$ for all $0<t\le t_\varepsilon$.
\item[$\RH$:] The Riemann zeta function $\zeta(s)=\sum_{n\ge1}n^{-s}$ ($\re s>1$) has a meromorphic continuation to $\C$ (Theorem~\ref{thm:FE} below), and every zero $\rho$ of this continuation with $0<\re\rho<1$ satisfies $\re\rho=\tfrac12$.
\end{itemize}
\end{definition}

\begin{theorem}[Main theorem]\label{thm:main}
$\GR\iff\RH$.
\end{theorem}

\begin{corollary}\label{cor:truth}
$\neg\GR\iff\neg\RH$.
A proof of either proposition yields, through Sections~\ref{sec:forward} and~\ref{sec:reverse}, a proof of the other; a counterexample to either yields a refutation of the other.
\end{corollary}
\begin{proof}
Immediate from Theorem~\ref{thm:main}, constructively: from $P\leftrightarrow Q$ one obtains $\neg P\leftrightarrow\neg Q$ by composing with the implications.
\end{proof}

\begin{remark}
$\GR$ is a statement about the quantitative Goldbach field $R(N)$ through its Laplace transform.
Remark~\ref{rem:finite} shows that the finite prefix $R(N)$, $N\le L$, determines the quantity at scale $t$ up to an error smaller than the critical rate once $L\gg t^{-1}\log(1/t)$, so $\GR$ is a statement about finite Goldbach data at every scale.
Nothing in this file concerns the proposition ``$R(N)>0$ for all even $N\ge 4$'', which is a different proposition.
\end{remark}

\section{The base}\label{sec:base}

Everything below is derived from the following, and from nothing else.

\begin{itemize}
\item[(B0)] \emph{Numbers and elementary functions.} The real and complex numbers; sequences and series; the complex exponential $\exp(z)=\sum z^{n}/n!$, which is complex differentiable with $\exp'=\exp$, satisfies $\exp(z+w)=\exp z\exp w$, $\exp z\ne0$, $|\exp z|=e^{\re z}$, and $\exp z=1$ iff $z\in2\pi i\Z$; $\cos$, $\sin$ defined by $\exp(i\theta)=\cos\theta+i\sin\theta$, so $\sin(\pi s)=0$ iff $s\in\Z$ and $e^{i\pi}=-1$; the real logarithm and real powers $x^{s}=\exp(s\log x)$ for $x>0$, with $|x^{s}|=x^{\re s}$; the inverse tangent, with $\arctan'(y)=1/(1+y^{2})$ and $\arctan(\pm1)=\pm\pi/4$; the binomial theorem and the geometric series. Connectedness: a locally constant function on a connected open subset of $\C$ is constant; an open disc, a convex open set and a punctured disc are connected.
\item[(B1)] \emph{Riemann integration.} For continuous functions on compact intervals, and for absolutely convergent improper integrals: linearity, the fundamental theorem of calculus, integration by parts, substitution, the triangle inequality $|\int f|\le\int|f|$. Uniform convergence: a uniform limit of continuous functions is continuous and may be integrated termwise on compact intervals; the Weierstrass $M$-test; a series of $C^{1}$ functions whose termwise derivative series converges uniformly on compacts may be differentiated termwise; a power series converges absolutely and uniformly on compact subsets of its disc of convergence and defines a continuous function there. Differentiation under the integral sign for an integrand with continuous partial derivative on a compact rectangle.
\item[(B2)] \emph{Lebesgue integration on $\R$ and $\R^{2}$.} Monotone and dominated convergence; Tonelli and Fubini; a function of two variables that is continuous in one and measurable in the other is jointly measurable; for continuous functions the Lebesgue integral agrees with the Riemann integral, absolutely convergent improper integrals included.
\end{itemize}

\begin{lemma}[Differentiation under the integral sign]\label{lem:param-diff}
Let $I\subset\R$ be an open interval, $(X,\mu)$ a measure space, and $F:I\times X\to\C$ such that $F(\cdot,x)$ is $C^{1}$ for each $x$, $F(s,\cdot)$ and $\partial_sF(s,\cdot)$ are measurable for each $s$, and $|F(s,x)|+|\partial_sF(s,x)|\le g(x)$ on $I\times X$ with $\int g\dd\mu<\infty$.
Then $f(s)=\int_XF(s,x)\dd\mu(x)$ is $C^{1}$ on $I$ with $f'(s)=\int_X\partial_sF(s,x)\dd\mu(x)$.
\end{lemma}
\begin{proof}
For $h\ne0$ small, $\bigl(F(s+h,x)-F(s,x)\bigr)/h=\int_0^1\partial_sF(s+\theta h,x)\dd\theta$ is bounded by $g(x)$ and tends to $\partial_sF(s,x)$ as $h\to0$; dominated convergence gives the derivative, and the same theorem gives its continuity.
\end{proof}

\section{Complex analysis}\label{sec:complex}

Throughout, $\Omega$ denotes an open subset of $\C$. A function $f:\Omega\to\C$ is \emph{holomorphic} if it is complex differentiable at every point of $\Omega$; no continuity of $f'$ is assumed.
A \emph{path} is a piecewise $C^{1}$ map $\gamma:[a,b]\to\C$, with $\int_\gamma f=\int_a^b f(\gamma(t))\gamma'(t)\dd t$ for continuous $f$ on the image, length $L(\gamma)=\int_a^b|\gamma'|$, and the estimate $|\int_\gamma f|\le L(\gamma)\max_\gamma|f|$.
If $F$ is holomorphic on an open set containing $\gamma$ and $F'=f$ is continuous, then $(F\circ\gamma)'=f(\gamma)\gamma'$ and (B1) gives $\int_\gamma f=F(\gamma(b))-F(\gamma(a))$; in particular $\int_\gamma f=0$ for closed $\gamma$.
For $n\ge2$ and $a\notin\gamma$ this applies to $f(w)=(w-a)^{-n}$ with $F(w)=(w-a)^{1-n}/(1-n)$ on $\C\setminus\{a\}$.
Segments $[z_0,z_1]$ are parametrised affinely; triangles and rectangles are traversed counterclockwise.

\begin{lemma}[Goursat]\label{lem:goursat}
Let $E\subset\Omega$ be finite, and let $f$ be continuous on $\Omega$ and holomorphic on $\Omega\setminus E$. Then $\int_{\partial T}f=0$ for every closed triangle $T\subset\Omega$.
\end{lemma}
\begin{proof}
First suppose $T\cap E=\emptyset$. Put $I=|\int_{\partial T}f|$, and let $d,P$ be the diameter and perimeter of $T$.
Joining the midpoints of the sides divides $T$ into four congruent triangles; the integrals over their boundaries sum to $\int_{\partial T}f$ (interior edges cancel), so one of them, $T_1$, has $|\int_{\partial T_1}f|\ge I/4$.
Iterating gives $T\supset T_1\supset T_2\supset\cdots$ with $|\int_{\partial T_n}f|\ge4^{-n}I$, diameter $2^{-n}d$ and perimeter $2^{-n}P$.
Let $z_0\in\bigcap_nT_n$. Since $f$ is differentiable at $z_0$, $f(z)=f(z_0)+f'(z_0)(z-z_0)+\eta(z)(z-z_0)$ with $\eta(z)\to0$ as $z\to z_0$.
The affine part has a primitive, so $\int_{\partial T_n}f=\int_{\partial T_n}\eta(z)(z-z_0)\dd z$, whence
$4^{-n}I\le 2^{-n}P\cdot2^{-n}d\cdot\sup_{T_n}|\eta|$, i.e.\ $I\le Pd\sup_{T_n}|\eta|\to0$.

If $T$ meets $E$, subdivide $T$ (by segments from the points of $E\cap T$) into finitely many triangles each containing at most one point of $E$, and that point as a vertex; it suffices to treat such a triangle $T'$ with exceptional vertex $p$.
Cutting off a small triangle $T''$ at $p$ by a segment parallel to the opposite side leaves pieces not meeting $E$, over which the integral vanishes, so $|\int_{\partial T'}f|=|\int_{\partial T''}f|\le L(\partial T'')\max_{T'}|f|$, which can be made arbitrarily small.
\end{proof}

\begin{lemma}[Primitives on convex sets; Cauchy's theorem]\label{lem:primitive}
Let $U$ be convex and open, and let $f$ be continuous on $U$ with $\int_{\partial T}f=0$ for every closed triangle $T\subset U$ (by Lemma~\ref{lem:goursat} this holds if $f$ is holomorphic on $U$ minus a finite set).
Fix $z_0\in U$ and put $F(z)=\int_{[z_0,z]}f$. Then $F$ is holomorphic on $U$ with $F'=f$, and $\int_\gamma f=0$ for every closed path $\gamma$ in $U$.
\end{lemma}
\begin{proof}
For $z,z+h\in U$ the triangle with vertices $z_0,z,z+h$ lies in $U$, so $F(z+h)-F(z)=\int_{[z,z+h]}f=hf(z)+\int_{[z,z+h]}\bigl(f(w)-f(z)\bigr)\dd w$, and the last term is at most $|h|\max_{w\in[z,z+h]}|f(w)-f(z)|=o(|h|)$ by continuity of $f$.
The second claim follows from the remark preceding Lemma~\ref{lem:goursat}, $F'=f$ being continuous.
\end{proof}

\begin{lemma}[Cauchy's integral formula on a disc; regularity]\label{lem:cauchy-disc}
Let $E\subset\Omega$ be finite, $f$ continuous on $\Omega$ and holomorphic on $\Omega\setminus E$, and let the closed disc $\overline D=\{|w-a|\le r\}$ lie in $\Omega$. Then for $|z-a|<r$,
\begin{equation}\label{eq:cif}
f(z)=\frac1{2\pi i}\oint_{|w-a|=r}\frac{f(w)}{w-z}\dd w .
\end{equation}
Consequently: $f$ is holomorphic on all of $\Omega$ (the points of $E$ are removable); $f$ has complex derivatives of all orders, with
$f^{(n)}(z)=\frac{n!}{2\pi i}\oint_{|w-a|=r}f(w)(w-z)^{-n-1}\dd w$; on the open disc $|z-a|<r$,
\begin{equation}\label{eq:taylor}
f(z)=\sum_{n\ge0}c_n(z-a)^{n},\qquad c_n=\frac{f^{(n)}(a)}{n!}=\frac1{2\pi i}\oint_{|w-a|=r}\frac{f(w)}{(w-a)^{n+1}}\dd w,
\end{equation}
the series converging absolutely and uniformly on compact subsets of the disc; and if $|f|\le M$ on $|w-a|=r$ then $|f^{(n)}(a)|\le n!\,M\,r^{-n}$ (Cauchy's estimates).
\end{lemma}
\begin{proof}
Since $\overline D$ is compact in $\Omega$, some open disc $U\supset\overline D$ lies in $\Omega$.
First let $z\notin E$, $|z-a|<r$, and define $g(w)=(f(w)-f(z))/(w-z)$ for $w\ne z$, $g(z)=f'(z)$.
Then $g$ is continuous on $U$ and holomorphic on $U\setminus(E\cup\{z\})$, so Lemma~\ref{lem:primitive} gives $\oint_{|w-a|=r}g=0$, i.e.\ $\oint f(w)/(w-z)\dd w=f(z)\oint\dd w/(w-z)$.
For $|z-a|<r=|w-a|$, $\frac1{w-z}=\sum_{n\ge0}\frac{(z-a)^{n}}{(w-a)^{n+1}}$ uniformly in $w$ on the circle, and with $w=a+re^{i\theta}$,
$\oint(w-a)^{-n-1}\dd w=\int_0^{2\pi}ir^{-n}e^{-in\theta}\dd\theta$ equals $2\pi i$ for $n=0$ and $0$ for $n\ge1$. Hence $\oint\dd w/(w-z)=2\pi i$ and \eqref{eq:cif} holds for $z\notin E$.
Both sides of \eqref{eq:cif} are continuous in $z$ on the open disc (the right side by (B1)), and $E$ is finite, so \eqref{eq:cif} holds for all $|z-a|<r$.

Write $F_n(z)=\oint f(w)(w-z)^{-n}\dd w$. For $w$ on the circle and $z,z+h$ in a smaller concentric disc, $\bigl((w-z-h)^{-n}-(w-z)^{-n}\bigr)/h\to n(w-z)^{-n-1}$ uniformly in $w$ as $h\to0$ (expand the difference of powers), so $F_n$ is complex differentiable with $F_n'=nF_{n+1}$.
Since $f=F_1/2\pi i$ on the disc, $f$ has derivatives of all orders there with the stated formula; as every point of $\Omega$ is the centre of such a disc, $f$ is holomorphic on $\Omega$.
Inserting the uniformly convergent expansion of $1/(w-z)$ into \eqref{eq:cif} and integrating termwise gives \eqref{eq:taylor}; $c_n=f^{(n)}(a)/n!$ by the derivative formula at $z=a$; absolute and uniform convergence on $|z-a|\le\rho<r$ follows from $|c_n|\le Mr^{-n}$, which is also Cauchy's estimate.
\end{proof}

\begin{corollary}[Morera; Weierstrass; parameter integrals]\label{cor:morera}
\begin{itemize}
\item[(a)] If $f$ is continuous on $\Omega$ and $\int_{\partial T}f=0$ for every closed triangle $T\subset\Omega$, then $f$ is holomorphic.
\item[(b)] If $f_n$ are holomorphic on $\Omega$ and $f_n\to f$ uniformly on compact subsets, then $f$ is holomorphic and $f_n^{(k)}\to f^{(k)}$ uniformly on compact subsets for every $k$. In particular a series of holomorphic functions converging uniformly on compacts is holomorphic and may be differentiated termwise.
\item[(c)] Let $(X,\mu)$ be a measure space and $F:\Omega\times X\to\C$ with $F(\cdot,x)$ holomorphic for each $x$ and $F(s,\cdot)$ measurable for each $s$. Suppose every $s_0\in\Omega$ has a neighbourhood $V$ and an integrable $g$ with $|F(s,x)|\le g(x)$ on $V\times X$. Then $f(s)=\int_XF(s,x)\dd\mu(x)$ is holomorphic on $\Omega$.
\end{itemize}
\end{corollary}
\begin{proof}
(a) On each open disc $D\subset\Omega$, Lemma~\ref{lem:primitive} gives a holomorphic $G$ with $G'=f$; by Lemma~\ref{lem:cauchy-disc} $G$ has derivatives of all orders, so $f=G'$ is holomorphic on $D$.
(b) $f$ is continuous, and $\int_{\partial T}f=\lim\int_{\partial T}f_n=0$, so $f$ is holomorphic by (a). For a closed disc $\{|w-a|\le r\}\subset\Omega$ and $|z-a|\le r/2$, the formula $f^{(k)}(z)-f_n^{(k)}(z)=\frac{k!}{2\pi i}\oint(f-f_n)(w)(w-z)^{-k-1}\dd w$ gives
\[
|f^{(k)}(z)-f_n^{(k)}(z)|\le k!\,2^{k+1}r^{-k}\max_{|w-a|=r}|f-f_n|\qquad(|z-a|\le r/2).
\]
(c) Continuity of $f$ at $s_0$: for $s_n\to s_0$ in $V$, $F(s_n,x)\to F(s_0,x)$ with domination $g$, so dominated convergence applies. For a closed triangle $T$ in a disc $V$ as in the hypothesis, $(t,x)\mapsto F(\gamma(t),x)\gamma'(t)$ (with $\gamma$ parametrising $\partial T$) is jointly measurable and dominated by $g(x)\max|\gamma'|$, so by Fubini $\int_{\partial T}f=\int_X\int_{\partial T}F(s,x)\dd s\dd\mu(x)=0$ by Lemma~\ref{lem:goursat}. Now apply (a).
\end{proof}

\begin{lemma}[Identity theorem; order of a zero]\label{lem:identity}
Let $\Omega$ be connected and $f$ holomorphic on $\Omega$.
\begin{itemize}
\item[(a)] If the zero set of $f$ has an accumulation point in $\Omega$ (for instance if $f$ vanishes on a nonempty open subset), then $f\equiv0$. Hence two holomorphic functions on $\Omega$ that agree on a nonempty open subset agree everywhere.
\item[(b)] If $f\not\equiv0$ and $f(a)=0$, there are a unique $m\ge1$ (the order of the zero) and a holomorphic $g$ on $\Omega$ with $f(z)=(z-a)^{m}g(z)$ and $g(a)\ne0$. Zeros of $f$ are isolated.
\end{itemize}
\end{lemma}
\begin{proof}
Let $Z=\{z\in\Omega:f^{(n)}(z)=0\text{ for all }n\ge0\}$. $Z$ is closed in $\Omega$ (each $f^{(n)}$ is continuous) and open (on a disc about a point of $Z$, \eqref{eq:taylor} gives $f\equiv0$).
Let $a\in\Omega$ be an accumulation point of zeros and suppose $a\notin Z$; let $m$ be least with $c_m=f^{(m)}(a)/m!\ne0$ in \eqref{eq:taylor} on a disc $D$ about $a$.
Then $g(z)=f(z)/(z-a)^{m}$ is holomorphic on $D\setminus\{a\}$ and $g(z)=\sum_{n\ge m}c_n(z-a)^{n-m}$ there, which by (B1) extends continuously to $a$ with value $c_m\ne0$; so $f=(z-a)^{m}g$ has no zero in a punctured neighbourhood of $a$, a contradiction. Thus $a\in Z$, $Z\ne\emptyset$, and connectedness gives $Z=\Omega$, i.e.\ $f\equiv0$.
For (b), with $m$ and $g$ as above, $g$ is continuous on $D$ and holomorphic on $D\setminus\{a\}$ and also on $\Omega\setminus\{a\}$, hence holomorphic on $\Omega$ by Lemma~\ref{lem:cauchy-disc}; uniqueness of $m$ and isolation of zeros follow from $g(a)\ne0$.
\end{proof}

\begin{lemma}[Maximum principle; Schwarz lemma]\label{lem:maxmod}
\begin{itemize}
\item[(a)] Let $\Omega$ be connected and $f$ holomorphic on $\Omega$. If $|f|$ attains a maximum at a point of $\Omega$, or if $\re f$ attains a maximum at a point of $\Omega$, then $f$ is constant.
\item[(b)] Let $f$ be holomorphic on an open set containing $\{|z|\le R\}$, with $f(0)=0$ and $|f|\le1$ on $|z|\le R$. Then $|f(z)|\le|z|/R$ for $|z|\le R$.
\end{itemize}
\end{lemma}
\begin{proof}
(a) Suppose $|f(z)|\le|f(a)|$ for all $z\in\Omega$, and let $\{|z-a|<R'\}\subset\Omega$. For $0<r<R'$, by \eqref{eq:taylor} the series $f(a+re^{i\theta})=\sum c_nr^{n}e^{in\theta}$ converges absolutely and uniformly in $\theta$, so the double series for $|f|^{2}=f\bar f$ may be integrated termwise; since $\int_0^{2\pi}e^{i(n-m)\theta}\dd\theta=2\pi\delta_{nm}$,
\[
\frac1{2\pi}\int_0^{2\pi}|f(a+re^{i\theta})|^{2}\dd\theta=\sum_{n\ge0}|c_n|^{2}r^{2n}\le|f(a)|^{2}=|c_0|^{2},
\]
so $c_n=0$ for $n\ge1$ and $f$ is constant on the disc, hence on $\Omega$ by Lemma~\ref{lem:identity}.
If $\re f\le\re f(a)$ on $\Omega$, apply this to $e^{f}$: $|e^{f}|=e^{\re f}$ attains a maximum at $a$, so $e^{f}$ is constant, $f'e^{f}=0$, $f'\equiv0$; then $f$ is constant on every disc in $\Omega$ (integrate $f'$ along segments) and hence on $\Omega$ by connectedness.

(b) $g(z)=f(z)/z$ for $z\ne0$, $g(0)=f'(0)$, is continuous and holomorphic off $0$, hence holomorphic by Lemma~\ref{lem:cauchy-disc}. On $|z|=R$, $|g|\le1/R$. If $\max_{|z|\le R}|g|$ were attained at an interior point, (a) on the disc $|z|<R$ would make $g$ constant there, and by continuity $|g|\equiv|g(z_1)|$ for some $|z_1|=R$; in all cases $\max_{|z|\le R}|g|\le1/R$.
\end{proof}

\begin{lemma}[Logarithms on convex sets]\label{lem:log}
Let $U$ be convex and open and $f$ holomorphic and nowhere zero on $U$. For $z_0\in U$ and $\ell_0$ with $e^{\ell_0}=f(z_0)$ there is a holomorphic $F$ on $U$ with $e^{F}=f$ and $F(z_0)=\ell_0$. Any two holomorphic logarithms of $f$ on $U$ differ by a constant in $2\pi i\Z$.
\end{lemma}
\begin{proof}
$f'/f$ is holomorphic on $U$ (Lemma~\ref{lem:cauchy-disc}), so by Lemma~\ref{lem:primitive} it has a primitive $F$, normalised by $F(z_0)=\ell_0$. Then $(fe^{-F})'=(f'-fF')e^{-F}=0$, so $fe^{-F}$ is constant, equal to $f(z_0)e^{-\ell_0}=1$.
If $e^{F_1}=e^{F_2}=f$ then $F_1-F_2$ is continuous with values in $2\pi i\Z$ (B0), hence constant on the connected set $U$.
\end{proof}

\begin{definition}[Poles, meromorphic functions, residues]\label{def:mero}
Let $S\subset\Omega$ be closed in $\Omega$ and discrete (no accumulation point in $\Omega$), and let $f$ be holomorphic on $\Omega\setminus S$.
A point $a\in S$ is \emph{removable} if $f$ extends holomorphically to a neighbourhood of $a$, and a \emph{pole of order $m\ge1$} if $(z-a)^{m}f(z)$ extends holomorphically to a neighbourhood of $a$ with nonzero value at $a$.
$f$ is \emph{meromorphic on $\Omega$} if every point of $S$ is removable or a pole; removable points are tacitly filled in, and the set of poles is still closed and discrete.
\end{definition}

\begin{lemma}[Local structure of meromorphic functions]\label{lem:mero}
\begin{itemize}
\item[(a)] If $f$ has a pole of order $m$ at $a$, there are unique $c_{-1},\dots,c_{-m}$, $c_{-m}\ne0$, and a holomorphic $h$ near $a$ with $f(z)=P_a(z)+h(z)$, $P_a(z)=\sum_{n=1}^{m}c_{-n}(z-a)^{-n}$. The \emph{principal part} $P_a$ is holomorphic on $\C\setminus\{a\}$ and $\res_af:=c_{-1}$ is the residue.
\item[(b)] If $f,g$ are holomorphic on connected $\Omega$ and $g\not\equiv0$, then $f/g$ is meromorphic on $\Omega$; its poles are among the zeros of $g$.
\item[(c)] If $h$ is meromorphic on connected $\Omega$ and $h\not\equiv0$, then $h'/h$ is meromorphic on $\Omega$, holomorphic away from the zeros and poles of $h$, with a simple pole of residue $m$ at every zero of order $m$ and a simple pole of residue $-m$ at every pole of order $m$.
\item[(d)] If $f$ has a pole at $a$ and $u$ is holomorphic near $a$ with $u(a)\ne0$, then $uf$ has a pole at $a$.
\item[(e)] If $\Omega$ is connected and $S\subset\Omega$ is closed and discrete, then $\Omega\setminus S$ is connected. Consequently two meromorphic functions on connected $\Omega$ that agree on a nonempty open subset have the same poles and agree everywhere else.
\end{itemize}
\end{lemma}
\begin{proof}
(a) Let $k(z)=(z-a)^{m}f(z)$, holomorphic near $a$ with $k(a)\ne0$, and expand $k=\sum_{j\ge0}k_j(z-a)^{j}$ by \eqref{eq:taylor}. Put $c_{-n}=k_{m-n}$ for $1\le n\le m$; then $f-P_a=(k(z)-\sum_{j<m}k_j(z-a)^{j})/(z-a)^{m}$, which is holomorphic off $a$ and, by (B1), extends continuously to $a$, hence is holomorphic near $a$ by Lemma~\ref{lem:cauchy-disc}. Uniqueness: two such decompositions give a principal part equal to a holomorphic function near $a$, forcing all its coefficients to vanish (multiply by $(z-a)^{m}$ and let $z\to a$ successively).
(b) The zeros of $g$ form a closed discrete set (Lemma~\ref{lem:identity}). At a zero $a$ of order $m$, $g=(z-a)^{m}g_1$ with $g_1(a)\ne0$, so $(z-a)^{m}f/g=f/g_1$ is holomorphic near $a$: if $f/g_1$ has a zero of order $k$ at $a$ (with $k=0$ if $(f/g_1)(a)\ne0$), then $f/g=(z-a)^{k-m}\cdot(\text{holomorphic, nonzero at }a)$, which is a pole of order $m-k$ if $k<m$ and removable otherwise.
(c) Near a zero $a$ of order $m$, $h=(z-a)^{m}g$ with $g(a)\ne0$ and $h'/h=m/(z-a)+g'/g$ with $g'/g$ holomorphic near $a$; near a pole of order $m$, $h=(z-a)^{-m}g$ and $h'/h=-m/(z-a)+g'/g$. Away from zeros and poles $h'/h$ is holomorphic. The zeros and poles of $h$ form a closed discrete set by Lemma~\ref{lem:identity}, applied on the connected set $\Omega\setminus\{\text{poles}\}$ (part (e)).
(d) $(z-a)^{m}uf$ extends holomorphically with value $u(a)\cdot\lim_{z\to a}(z-a)^{m}f(z)\ne0$.
(e) Suppose $\Omega\setminus S=V_1\sqcup V_2$ with $V_1,V_2$ open and nonempty. Each $a\in S$ has a punctured disc contained in $\Omega\setminus S$; being connected, it lies in $V_1$ or in $V_2$; assign $a$ to that side. The enlarged sets $V_1',V_2'$ are open, disjoint, nonempty, and cover $\Omega$, contradicting connectedness. For the consequence, apply Lemma~\ref{lem:identity} to the difference on the connected set $\Omega$ minus the union of the two pole sets; the difference is then $0$ there, so at each would-be pole it is removable with value $0$.
\end{proof}

\begin{lemma}[Cauchy's theorem and formula for rectangles]\label{lem:rect}
Let $Q=[x_0,x_1]\times[y_0,y_1]\subset\C$ be a closed rectangle and $f$ holomorphic on an open set $\Omega\supset Q$.
\begin{itemize}
\item[(a)] $\oint_{\partial Q}f=0$.
\item[(b)] For every $z$ in the interior of $Q$, $\oint_{\partial Q}\dfrac{f(w)}{w-z}\dd w=2\pi i\,f(z)$; in particular $\oint_{\partial Q}\dfrac{\dd w}{w-z}=2\pi i$.
\end{itemize}
\end{lemma}
\begin{proof}
(a) Every point of the compact set $Q$ has an open disc about it inside $\Omega$; by compactness there is $\delta>0$ such that every subset of $Q$ of diameter $<\delta$ lies in one such disc. Subdivide $Q$ into congruent subrectangles of diameter $<\delta$; each boundary integral vanishes by Lemma~\ref{lem:primitive} (a disc is convex), and interior edges cancel in the sum, which equals $\oint_{\partial Q}f$.

(b) Fix $z$ in the interior and $\rho>0$ so small that the square $S_\rho=z+[-\rho,\rho]^{2}$ lies in the interior of $Q$.
The two horizontal and two vertical lines through the sides of $S_\rho$ cut the closure of $Q\setminus S_\rho$ into eight closed rectangles $Q_1,\dots,Q_8$, none containing $z$. The function $w\mapsto f(w)/(w-z)$ is holomorphic on the open set $\Omega\setminus\{z\}\supset Q_j$, so by (a) its integral over each $\partial Q_j$ vanishes; summing, interior edges cancel and
\[
\oint_{\partial Q}\frac{f(w)}{w-z}\dd w=\oint_{\partial S_\rho}\frac{f(w)}{w-z}\dd w
=\oint_{\partial S_\rho}\frac{f(w)-f(z)}{w-z}\dd w+f(z)\oint_{\partial S_\rho}\frac{\dd w}{w-z}.
\]
The first integrand is bounded on a neighbourhood of $z$ (it extends continuously by $f'(z)$) and $L(\partial S_\rho)=8\rho$, so the first integral tends to $0$ as $\rho\to0$; the second is independent of $\rho$ by the substitution $w=z+\rho w'$, which reduces it to the integral over the boundary of the unit square $[-1,1]^{2}$ of $\dd w'/w'$.
Since the left side does not depend on $\rho$, it equals $f(z)\oint_{\partial[-1,1]^{2}}\dd w'/w'$.
On the right side $w'=1+iy$, $-1\le y\le1$:
\[
\int_{-1}^{1}\frac{i\dd y}{1+iy}=\int_{-1}^{1}\frac{y+i}{1+y^{2}}\dd y=i\bigl(\arctan1-\arctan(-1)\bigr)=\frac{i\pi}{2},
\]
the real part vanishing by oddness. Multiplication by $i$ maps the unit square onto itself, carries each side onto the next in the counterclockwise order, and leaves $\dd w'/w'$ invariant; so each side contributes $i\pi/2$ and the total is $2\pi i$.
\end{proof}

\begin{lemma}[Residue theorem for rectangles]\label{lem:residue}
Let $f$ be meromorphic on an open set $\Omega$ containing the closed rectangle $Q$, with no pole on $\partial Q$ and finitely many poles $a_1,\dots,a_k$ in the interior of $Q$. Then
$\oint_{\partial Q}f=2\pi i\sum_{j=1}^{k}\res_{a_j}f$.
\end{lemma}
\begin{proof}
Let $P_j$ be the principal part at $a_j$. By Lemma~\ref{lem:mero}(a), $g=f-\sum_jP_j$ is holomorphic on $\Omega$ (each $a_j$ is removable for $f-P_j$, and the other $P_i$ are holomorphic there), so $\oint_{\partial Q}g=0$ by Lemma~\ref{lem:rect}(a).
For $n\ge2$ the term $(w-a_j)^{-n}$ has a primitive on $\C\setminus\{a_j\}$, so its integral over the closed path $\partial Q$ vanishes; and $\oint_{\partial Q}\dd w/(w-a_j)=2\pi i$ by Lemma~\ref{lem:rect}(b).
\end{proof}

\section{Arithmetic}\label{sec:arith}

\begin{lemma}[B\'ezout; Euclid]\label{lem:euclid}
\begin{itemize}
\item[(a)] For integers $a,b$ not both $0$, the least positive element $d$ of $\{ax+by:x,y\in\Z\}$ divides $a$ and $b$, and every common divisor of $a$ and $b$ divides $d$.
\item[(b)] If $p$ is prime and $p\mid ab$ then $p\mid a$ or $p\mid b$.
\end{itemize}
\end{lemma}
\begin{proof}
(a) Write $a=qd+r$ with $0\le r<d$; then $r=a-q(ax+by)$ is in the set, so $r=0$ by minimality; likewise $d\mid b$. A common divisor of $a,b$ divides $ax+by=d$.
(b) If $p\nmid a$, the $d$ of (a) for $(p,a)$ is a positive divisor of $p$ other than $p$, so $d=1=px+ay$, and $b=pbx+aby$ is divisible by $p$.
\end{proof}

\begin{theorem}[Unique factorisation]\label{thm:uf}
Every integer $n\ge1$ is a product of primes, uniquely up to order. Writing $v_p(n)$ for the number of factors $p$: $a\mid b$ iff $v_p(a)\le v_p(b)$ for every prime $p$; and for positive integers $a_1,\dots,a_r$, $\lcm(a_1,\dots,a_r)=\prod_p p^{\max_iv_p(a_i)}$.
\end{theorem}
\begin{proof}
Existence by induction: $n=1$ is the empty product; otherwise the least divisor $p>1$ of $n$ is prime and $n/p<n$.
Uniqueness by induction on $n$: if $p_1\cdots p_r=q_1\cdots q_s$, Lemma~\ref{lem:euclid}(b) gives $p_1\mid q_j$ for some $j$, so $p_1=q_j$; cancel and apply the induction hypothesis.
If $a\mid b$ then $b=ac$ and $v_p(b)=v_p(a)+v_p(c)$ by uniqueness; conversely $\prod p^{v_p(b)-v_p(a)}$ is an integer $c$ with $ac=b$.
Hence $m$ is a common multiple of the $a_i$ iff $v_p(m)\ge\max_iv_p(a_i)$ for all $p$, and the least such $m$ is the displayed product.
\end{proof}

\section{The Gamma function}\label{sec:gamma}

\begin{definition}
For $\re s>0$ let $\Gamma(s)=\int_0^\infty e^{-x}x^{s-1}\dd x$.
\end{definition}

\begin{lemma}\label{lem:gamma-basic}
$\Gamma$ is holomorphic on $\re s>0$, $\Gamma(s+1)=s\Gamma(s)$, $\Gamma(1)=1$, $\Gamma(n)=(n-1)!$ for $n\ge1$.
$\Gamma$ extends to a meromorphic function on $\C$ whose poles are exactly the simple poles at $0,-1,-2,\dots$.
\end{lemma}
\begin{proof}
Holomorphy by Corollary~\ref{cor:morera}(c) with $g(x)=e^{-x}(x^{a-1}+x^{b-1})$ on $a\le\re s\le b$.
Integration by parts gives $\Gamma(s+1)=s\Gamma(s)$, and $\Gamma(1)=1$.
For $\re s>-n$ define $\Gamma(s)=\Gamma(s+n)/\bigl(s(s+1)\cdots(s+n-1)\bigr)$; the definitions agree on overlaps by the functional equation.
The only singularities are at $s=-j$, $j\ge0$: choosing $n>j$, the numerator there is $\Gamma(n-j)=(n-j-1)!\ne0$ while the denominator has a simple zero, so each is a simple pole.
\end{proof}

\begin{lemma}[Beta integral]\label{lem:beta}
For $\re a>0$, $\re b>0$:
$\displaystyle \Gamma(a)\Gamma(b)=\Gamma(a+b)\int_0^1 v^{a-1}(1-v)^{b-1}\dd v$.
\end{lemma}
\begin{proof}
$\Gamma(a)\Gamma(b)=\iint_{(0,\infty)^2}e^{-x-y}x^{a-1}y^{b-1}\dd x\dd y$ (absolutely convergent, Fubini).
Substitute $x=wv$, $y=w(1-v)$ with $w>0$, $0<v<1$; the Jacobian is $w$, and the integral becomes
$\int_0^\infty e^{-w}w^{a+b-1}\dd w\int_0^1 v^{a-1}(1-v)^{b-1}\dd v$.
\end{proof}

\begin{lemma}[Reflection]\label{lem:reflection}
For all $s\in\C\setminus\Z$: $\Gamma(s)\Gamma(1-s)=\pi/\sin(\pi s)$. In particular $\Gamma(\tfrac12)=\sqrt\pi$.
\end{lemma}
\begin{proof}
First let $0<\re s<1$. By Lemma~\ref{lem:beta} with $a=s$, $b=1-s$,
$\Gamma(s)\Gamma(1-s)=\int_0^1 v^{s-1}(1-v)^{-s}\dd v$.
Substituting $v=u/(1+u)$, $\dd v=\dd u/(1+u)^2$, $v^{s-1}(1-v)^{-s}=u^{s-1}(1+u)$, this equals
$I(s)=\int_0^\infty \frac{u^{s-1}}{1+u}\dd u=\int_{\R}\frac{e^{sv}}{1+e^{v}}\dd v$ (with $u=e^{v}$).
Let $g(z)=e^{sz}/(1+e^{z})$, meromorphic on $\C$ with poles where $e^{z}=-1$, i.e.\ at $z=i\pi(2k+1)$, and integrate over the rectangle with vertices $\pm T$, $\pm T+2\pi i$ (Lemma~\ref{lem:residue}; no pole lies on its boundary).
Inside there is exactly one pole, $z=i\pi$, where $1+e^{z}$ vanishes with derivative $e^{i\pi}=-1$; so $g=e^{sz}/\bigl((z-i\pi)u(z)\bigr)$ with $u$ holomorphic, $u(i\pi)=-1$, and $\res_{i\pi}g=-e^{i\pi s}$.
The top edge contributes $-\int_{-T}^{T}g(v+2\pi i)\dd v=-e^{2\pi i s}\int_{-T}^{T}g(v)\dd v$.
On the right edge $|g(T+iy)|\le e^{\sigma T}/(e^{T}-1)\to0$ since $\sigma=\re s<1$; on the left edge $|g(-T+iy)|\le e^{-\sigma T}/(1-e^{-T})\to0$ since $\sigma>0$.
Letting $T\to\infty$: $(1-e^{2\pi i s})I(s)=2\pi i\,(-e^{i\pi s})$, so
$I(s)=\dfrac{2\pi i}{e^{i\pi s}-e^{-i\pi s}}=\dfrac{\pi}{\sin\pi s}$.
Both sides of the claimed identity are meromorphic on $\C$ and agree on the strip, so they agree on $\C\setminus\Z$ by Lemma~\ref{lem:mero}(e).
At $s=\tfrac12$: $\Gamma(\tfrac12)^{2}=\pi$ and $\Gamma(\tfrac12)>0$.
\end{proof}

\begin{corollary}\label{cor:gamma-nonzero}
$\Gamma(s)\ne0$ for every $s\in\C$ that is not a pole; in particular $\Gamma(s)\ne 0$ for $\re s>0$, and $1/\Gamma$ is entire.
\end{corollary}
\begin{proof}
If $s\notin\Z$ then $\Gamma(s)\Gamma(1-s)=\pi/\sin\pi s$ is finite and nonzero, so $\Gamma(s)\ne0$.
If $s=n\ge1$ then $\Gamma(n)=(n-1)!\ne0$.
Nonpositive integers are poles.
Hence $1/\Gamma$ has no poles, and its zeros are the poles of $\Gamma$.
\end{proof}

\begin{lemma}[Decay on vertical strips]\label{lem:gamma-decay}
For $0<a\le b<\infty$ there is $K=K(a,b)$ with $|\Gamma(\sigma+i\tau)|\le K/(1+\tau^{2})$ for all $a\le\sigma\le b$, $\tau\in\R$.
\end{lemma}
\begin{proof}
Substituting $x=e^{u}$, $\Gamma(\sigma+i\tau)=\int_\R\varphi_\sigma(u)e^{i\tau u}\dd u$ with $\varphi_\sigma(u)=\exp(\sigma u-e^{u})$, a smooth function all of whose derivatives are integrable (they are polynomials in $e^{u}$ times $\varphi_\sigma$).
Two integrations by parts give, for $\tau\neq0$, $\Gamma(\sigma+i\tau)=-\tau^{-2}\int_\R\varphi_\sigma''(u)e^{i\tau u}\dd u$, and
$\varphi_\sigma''=\bigl((\sigma-e^{u})^{2}-e^{u}\bigr)\varphi_\sigma$, so
\[
|\Gamma(\sigma+i\tau)|\le\tau^{-2}\int_\R\bigl((\sigma-e^{u})^{2}+e^{u}\bigr)e^{\sigma u-e^{u}}\dd u
=\tau^{-2}\bigl(\Gamma(\sigma+2)-(2\sigma-1)\Gamma(\sigma+1)+\sigma^{2}\Gamma(\sigma)\bigr),
\]
using $\int_\R e^{ku}e^{\sigma u-e^{u}}\dd u=\Gamma(\sigma+k)$.
The right side is a continuous function of $\sigma\in[a,b]$, bounded by some $K_1$; also $|\Gamma(\sigma+i\tau)|\le\Gamma(\sigma)\le K_2$ on $[a,b]$.
Take $K=2(K_1+K_2)$.
\end{proof}

\section{Fourier analysis}\label{sec:fourier}

For an integrable $f:\R\to\C$ put $\widehat f(\xi)=\int_\R f(u)e^{-2\pi iu\xi}\dd u$.

\begin{lemma}[The Gaussian]\label{lem:gauss}
$\int_\R e^{-\pi u^{2}}\dd u=1$, and for $x>0$, $\xi\in\R$:
$\displaystyle\int_\R e^{-\pi xu^{2}}e^{-2\pi iu\xi}\dd u=x^{-1/2}e^{-\pi\xi^{2}/x}$.
\end{lemma}
\begin{proof}
Substituting $x=\pi u^{2}$ in $\Gamma(\tfrac12)=\int_0^\infty e^{-x}x^{-1/2}\dd x$ gives
$\Gamma(\tfrac12)=2\sqrt\pi\int_0^\infty e^{-\pi u^{2}}\dd u$, so
$\int_\R e^{-\pi u^{2}}\dd u=\Gamma(\tfrac12)/\sqrt\pi=1$ by Lemma~\ref{lem:reflection}.
Let $h(\xi)=\int_\R e^{-\pi u^{2}}e^{-2\pi iu\xi}\dd u$. By Lemma~\ref{lem:param-diff} (domination by $2\pi|u|e^{-\pi u^{2}}$) and an integration by parts on $[-A,A]$ with $A\to\infty$,
\[
h'(\xi)=\int_\R(-2\pi iu)e^{-\pi u^{2}}e^{-2\pi iu\xi}\dd u
=i\int_\R\bigl(e^{-\pi u^{2}}\bigr)'e^{-2\pi iu\xi}\dd u
=i\int_\R e^{-\pi u^{2}}(2\pi i\xi)e^{-2\pi iu\xi}\dd u
=-2\pi\xi\,h(\xi),
\]
so $(h(\xi)e^{\pi\xi^{2}})'=0$ and $h(\xi)=h(0)e^{-\pi\xi^{2}}=e^{-\pi\xi^{2}}$.
For general $x$ substitute $u=v/\sqrt x$.
\end{proof}

\begin{lemma}[Fourier inversion]\label{lem:inversion}
Let $f:\R\to\C$ be continuous, bounded and integrable, with $\widehat f$ integrable. Then $f(u)=\int_\R\widehat f(\xi)e^{2\pi iu\xi}\dd\xi$ for every $u\in\R$.
\end{lemma}
\begin{proof}
For $\delta>0$ put $I_\delta(u)=\int_\R\widehat f(\xi)e^{-\pi\delta^{2}\xi^{2}}e^{2\pi iu\xi}\dd\xi$. Since $|f(y)|e^{-\pi\delta^{2}\xi^{2}}$ is integrable on $\R^{2}$, Fubini and Lemma~\ref{lem:gauss} (with $x=\delta^{2}$) give
\[
I_\delta(u)=\int_\R f(y)\Bigl(\int_\R e^{-\pi\delta^{2}\xi^{2}}e^{-2\pi i(y-u)\xi}\dd\xi\Bigr)\dd y
=\int_\R f(y)\,\delta^{-1}e^{-\pi(y-u)^{2}/\delta^{2}}\dd y=\int_\R f(u+\delta v)e^{-\pi v^{2}}\dd v .
\]
As $\delta\to0$ the last integrand tends to $f(u)e^{-\pi v^{2}}$ pointwise and is dominated by $\sup|f|\,e^{-\pi v^{2}}$, so $I_\delta(u)\to f(u)$; and $I_\delta(u)\to\int\widehat f(\xi)e^{2\pi iu\xi}\dd\xi$ by dominated convergence with dominant $|\widehat f|$.
\end{proof}

\begin{lemma}[Fourier series of $C^{2}$ periodic functions]\label{lem:fourier-series}
Let $F:\R\to\C$ be $C^{2}$ and $1$-periodic, and $c_k=\int_0^1F(x)e^{-2\pi ikx}\dd x$. Then $\sum_k|c_k|<\infty$ and $F(x)=\sum_{k\in\Z}c_ke^{2\pi ikx}$ for every $x$.
\end{lemma}
\begin{proof}
Two integrations by parts (boundary terms cancel by periodicity) give $c_k=-(2\pi k)^{-2}\int_0^1F''(x)e^{-2\pi ikx}\dd x$, so $|c_k|\le\max|F''|/(4\pi^{2}k^{2})$ for $k\ne0$.
Let $S_N(x)=\sum_{|k|\le N}c_ke^{2\pi ikx}=\int_{-1/2}^{1/2}F(x-y)D_N(y)\dd y$ with $D_N(y)=\sum_{|k|\le N}e^{2\pi iky}$ (substitute $y=x-x'$ and use periodicity).
Summing the geometric progression, $D_N(y)=\sin\bigl((2N+1)\pi y\bigr)/\sin(\pi y)$ for $y\notin\Z$, and $\int_{-1/2}^{1/2}D_N=1$.
Hence $S_N(x)-F(x)=\int_{-1/2}^{1/2}g_x(y)\sin\bigl((2N+1)\pi y\bigr)\dd y$ with $g_x(y)=(F(x-y)-F(x))/\sin(\pi y)$.
Since $F(x-y)-F(x)=-y\int_0^1F'(x-ty)\dd t$ and $y/\sin(\pi y)$ extends to a $C^{1}$ function on $[-\tfrac12,\tfrac12]$, $g_x$ extends to a $C^{1}$ function there; one integration by parts gives
\[
|S_N(x)-F(x)|\le\frac{|g_x(\tfrac12)|+|g_x(-\tfrac12)|+\int_{-1/2}^{1/2}|g_x'|}{(2N+1)\pi}\longrightarrow0 .\qedhere
\]
\end{proof}

\begin{lemma}[Poisson summation]\label{lem:poisson}
Let $f\in C^{2}(\R)$ with $|f(u)|,|f'(u)|,|f''(u)|\le C/(1+u^{2})$ for all $u$. Then $\sum_{n\in\Z}f(n)=\sum_{n\in\Z}\widehat f(n)$, both series converging absolutely.
\end{lemma}
\begin{proof}
For $x\in[-1,2]$ and $|n|\ge3$, $1+(x+n)^{2}\ge1+(|n|-2)^{2}$, so the series $\sum_nf(x+n)$, $\sum_nf'(x+n)$, $\sum_nf''(x+n)$ converge uniformly on $[-1,2]$, and likewise on every compact interval. By (B1), $F(x)=\sum_nf(x+n)$ is $C^{2}$ on $\R$, and it is $1$-periodic.
Its Fourier coefficients are, by termwise integration of the uniformly convergent series on $[0,1]$,
$c_k=\sum_n\int_0^1f(x+n)e^{-2\pi ikx}\dd x=\sum_n\int_n^{n+1}f(y)e^{-2\pi iky}\dd y=\widehat f(k)$.
Lemma~\ref{lem:fourier-series} at $x=0$ gives $\sum_nf(n)=F(0)=\sum_kc_k=\sum_k\widehat f(k)$.
\end{proof}

\begin{corollary}[Cahen--Mellin integral]\label{lem:cahen-mellin}
For $c>0$ and $x>0$:
$\displaystyle e^{-x}=\frac{1}{2\pi}\int_\R\Gamma(c+i\tau)\,x^{-c-i\tau}\dd\tau$.
\end{corollary}
\begin{proof}
With $\varphi_c(u)=\exp(cu-e^{u})$, continuous, bounded and integrable,
$\widehat{\varphi_c}(\xi)=\int_\R e^{cu-e^{u}}e^{-2\pi i\xi u}\dd u=\Gamma(c-2\pi i\xi)$, which is integrable in $\xi$ by Lemma~\ref{lem:gamma-decay}.
Lemma~\ref{lem:inversion} gives, with $\tau=-2\pi\xi$,
\[
\varphi_c(u)=\int_\R\Gamma(c-2\pi i\xi)e^{2\pi iu\xi}\dd\xi=\frac1{2\pi}\int_\R\Gamma(c+i\tau)e^{-i\tau u}\dd\tau .
\]
Put $u=\log x$: $x^{c}e^{-x}=\frac1{2\pi}\int_\R\Gamma(c+i\tau)x^{-i\tau}\dd\tau$.
\end{proof}

\section{The zeta function}\label{sec:zeta}

\begin{lemma}[Euler product and logarithm]\label{lem:euler}
For $\re s>1$ the series $L(s)=\sum_{p}\sum_{k\ge1}\frac{p^{-ks}}{k}=\sum_{n\ge2}\frac{\Lam(n)}{\log n}n^{-s}$ converges absolutely and locally uniformly, $\zeta(s)=\exp L(s)$, hence $\zeta(s)\neq0$, and
\[
-\frac{\zeta'}{\zeta}(s)=\sum_{n\ge1}\Lam(n)n^{-s}\qquad(\re s>1).
\]
Moreover $|L(2+i\tau)|\le L(2)=\log\zeta(2)<1$ for all real $\tau$.
\end{lemma}
\begin{proof}
The two expressions for $L$ agree termwise: for $n=p^{k}$, $1/k=\log p/\log n=\Lam(n)/\log n$.
Absolute convergence: $\sum_p\sum_k|p^{-ks}|/k\le\sum_{n\ge2}n^{-\sigma}<\infty$ for $\sigma=\re s>1$, locally uniformly, so $L$ is holomorphic (Corollary~\ref{cor:morera}(b)).
For a fixed prime $p$, $\exp\bigl(\sum_k p^{-ks}/k\bigr)=(1-p^{-s})^{-1}$ by the power series of $-\log(1-z)$, $|z|<1$.
Hence for a bound $P$, $\exp\bigl(\sum_{p\le P}\sum_k p^{-ks}/k\bigr)=\prod_{p\le P}\sum_{k\ge0}p^{-ks}=\sum_{n\in S_P}n^{-s}$, where $S_P$ is the set of positive integers all of whose prime factors are at most $P$; the last equality is Theorem~\ref{thm:uf}, the expansion of the finite product of absolutely convergent series.
As $P\to\infty$ the left side tends to $\exp L(s)$ and the right side to $\zeta(s)$ (the omitted terms are bounded by $\sum_{n>P}n^{-\sigma}\to0$).
Termwise differentiation (Corollary~\ref{cor:morera}(b)): $L'(s)=-\sum_p\sum_k(\log p)p^{-ks}=-\sum_{n}\Lam(n)n^{-s}$, and $\zeta'/\zeta=L'$.
Finally $|L(2+i\tau)|\le\sum_p\sum_k p^{-2k}/k=L(2)=\log\zeta(2)$ and $\zeta(2)\le 1+\int_1^\infty x^{-2}\dd x=2<e$.
\end{proof}

\begin{lemma}[Continuation to $\re s>0$ and a growth bound]\label{lem:continuation}
For every integer $N\ge1$ and every $s$ with $\re s>0$, $s\ne1$,
\begin{equation}\label{eq:partial}
\zeta(s)=\sum_{n\le N}n^{-s}+\frac{N^{1-s}}{s-1}-s\int_N^\infty\{x\}x^{-s-1}\dd x,
\end{equation}
where $\{x\}=x-\lfloor x\rfloor$; the right side defines a meromorphic continuation of $\zeta$ to $\re s>0$ whose only pole is simple, at $s=1$, with residue $1$.
Consequently $-\zeta'/\zeta$ is meromorphic on $\re s>0$, with a simple pole of residue $1$ at $s=1$ and simple poles of residue $-m$ at each zero of $\zeta$ of multiplicity $m$, and no other poles.
Moreover
\begin{equation}\label{eq:growth}
|\zeta(\sigma+i\tau)|\le 12\sqrt{|\tau|}\qquad\text{for }\tfrac12\le\sigma\le4,\ |\tau|\ge2.
\end{equation}
\end{lemma}
\begin{proof}
For integers $M>N$ and $n\le x<n+1$ one has $\lfloor x\rfloor=n$, so $s\int_n^{n+1}\lfloor x\rfloor x^{-s-1}\dd x=n\bigl(n^{-s}-(n+1)^{-s}\bigr)$; summing over $N\le n<M$ and rearranging,
$\sum_{N<n\le M}n^{-s}=M^{1-s}-N^{1-s}+s\int_N^M\lfloor x\rfloor x^{-s-1}\dd x$.
For $\re s>1$ let $M\to\infty$ and write $\lfloor x\rfloor=x-\{x\}$, $s\int_N^\infty x^{-s}\dd x=sN^{1-s}/(s-1)$; this gives \eqref{eq:partial}.
The integral converges absolutely for $\re s>0$, dominated locally uniformly by $\int_N^\infty x^{-\sigma_0-1}\dd x$, hence is holomorphic there (Corollary~\ref{cor:morera}(c)); the continuation is unique by Lemma~\ref{lem:mero}(e).
With $N=1$, $\zeta(s)=\frac{s}{s-1}-s\int_1^\infty\{x\}x^{-s-1}\dd x$, so $(s-1)\zeta(s)\to1$ as $s\to1$.
The statement about $-\zeta'/\zeta$ is Lemma~\ref{lem:mero}(c).

For \eqref{eq:growth}: take $N=\lfloor|\tau|\rfloor$, so $|\tau|/2\le N\le|\tau|$.
Then $|\sum_{n\le N}n^{-s}|\le\sum_{n\le N}n^{-1/2}\le 2\sqrt N\le2\sqrt{|\tau|}$;
$|N^{1-s}/(s-1)|\le N^{1/2}/|\tau|\le 1$;
and $|s\int_N^\infty\{x\}x^{-s-1}\dd x|\le|s|N^{-\sigma}/\sigma\le(4+|\tau|)\cdot2\cdot\sqrt{2/|\tau|}\le 9\sqrt{|\tau|}$ for $|\tau|\ge 2$.
\end{proof}

\begin{lemma}[No zeros on $\re s=1$]\label{lem:mertens}
$\zeta(1+i\tau)\ne0$ for every real $\tau\ne0$. Together with Lemma~\ref{lem:euler}, $\zeta(s)\ne0$ for $\re s\ge1$, $s\ne1$.
\end{lemma}
\begin{proof}
For $\sigma>1$ and real $\tau$, by Lemma~\ref{lem:euler},
$\re L(\sigma+i\tau)=\sum_{n\ge2}\frac{\Lam(n)}{\log n}n^{-\sigma}\cos(\tau\log n)$.
Since $3+4\cos\theta+\cos2\theta=2(1+\cos\theta)^{2}\ge0$,
\[
3L(\sigma)+4\re L(\sigma+i\tau)+\re L(\sigma+2i\tau)\ge0,\qquad\text{so}\qquad
\zeta(\sigma)^{3}\,|\zeta(\sigma+i\tau)|^{4}\,|\zeta(\sigma+2i\tau)|\ge1 .
\]
Suppose $\zeta(1+i\tau_0)=0$ with $\tau_0\ne0$.
By Lemma~\ref{lem:continuation}, $\zeta$ is holomorphic on a neighbourhood of the segments $\{\sigma+i\tau_0\}$ and $\{\sigma+2i\tau_0\}$, $1\le\sigma\le2$, so there are constants with
$|\zeta(\sigma+i\tau_0)|\le K(\sigma-1)$ and $|\zeta(\sigma+2i\tau_0)|\le K'$ for $1<\sigma\le2$,
while $\zeta(\sigma)=\sigma/(\sigma-1)-\sigma\int_1^\infty\{x\}x^{-\sigma-1}\dd x\le\sigma/(\sigma-1)\le 2/(\sigma-1)$ by \eqref{eq:partial} with $N=1$, the subtracted integral being nonnegative for real $\sigma>1$.
Then the product is at most $8K^{4}K'(\sigma-1)\to0$ as $\sigma\downarrow1$, a contradiction.
\end{proof}

\begin{lemma}[Theta inversion]\label{lem:theta}
For $x>0$ let $\theta(x)=\sum_{n\in\Z}e^{-\pi n^{2}x}$ and $\omega(x)=\sum_{n\ge1}e^{-\pi n^{2}x}$, so $\theta=1+2\omega$.
Then $\theta(1/x)=\sqrt{x}\,\theta(x)$, and $0<\omega(x)\le 2e^{-\pi x}$ for $x\ge1$.
\end{lemma}
\begin{proof}
The function $f(u)=e^{-\pi xu^{2}}$ satisfies the hypotheses of Lemma~\ref{lem:poisson}, and $\widehat f(\xi)=x^{-1/2}e^{-\pi\xi^{2}/x}$ by Lemma~\ref{lem:gauss}; Poisson summation gives $\theta(x)=x^{-1/2}\theta(1/x)$.
For $x\ge1$: $\omega(x)\le\sum_{n\ge1}e^{-\pi nx}=e^{-\pi x}/(1-e^{-\pi x})\le2e^{-\pi x}$.
\end{proof}

\begin{theorem}[Functional equation]\label{thm:FE}
Let $\xi(s)=\pi^{-s/2}\Gamma(s/2)\zeta(s)$ for $\re s>1$. Then
\begin{equation}\label{eq:xi}
\xi(s)=\frac{1}{s(s-1)}+\int_1^\infty\omega(x)\bigl(x^{s/2-1}+x^{(1-s)/2-1}\bigr)\dd x ,
\end{equation}
where the integral is an entire function of $s$.
Consequently $\zeta$ extends to a meromorphic function on $\C$, agreeing with Lemma~\ref{lem:continuation} on $\re s>0$, and
\[
\pi^{-s/2}\Gamma(s/2)\zeta(s)=\pi^{-(1-s)/2}\Gamma\bigl(\tfrac{1-s}{2}\bigr)\zeta(1-s)
\]
as meromorphic functions on $\C$.
\end{theorem}
\begin{proof}
For $\re s>1$ and $n\ge1$, substituting $y=\pi n^{2}x$ gives $\int_0^\infty e^{-\pi n^{2}x}x^{s/2-1}\dd x=(\pi n^{2})^{-s/2}\Gamma(s/2)$.
Summing over $n$ and interchanging, which is legitimate because
\[
\sum_{n\ge1}\int_0^\infty e^{-\pi n^2x}x^{\sigma/2-1}\dd x=\pi^{-\sigma/2}\Gamma(\sigma/2)\zeta(\sigma)<\infty ,
\qquad\text{gives}\qquad
\xi(s)=\int_0^\infty\omega(x)x^{s/2-1}\dd x .
\]
Split at $x=1$. In $\int_0^1$ substitute $x=1/y$ and use $\omega(1/y)=-\tfrac12+\tfrac{\sqrt y}{2}+\sqrt y\,\omega(y)$ (from Lemma~\ref{lem:theta}):
\[
\int_0^1\omega(x)x^{s/2-1}\dd x=\int_1^\infty\omega(1/y)\,y^{-s/2-1}\dd y
=-\frac1s+\frac1{s-1}+\int_1^\infty\omega(y)y^{(1-s)/2-1}\dd y ,
\]
using $\int_1^\infty y^{-s/2-1}\dd y=2/s$ and $\int_1^\infty y^{-s/2-1/2}\dd y=2/(s-1)$ for $\re s>1$.
Since $-1/s+1/(s-1)=1/(s(s-1))$, \eqref{eq:xi} follows.
The integral in \eqref{eq:xi} is entire by Corollary~\ref{cor:morera}(c), since $\omega(x)\le2e^{-\pi x}$ dominates $x^{\sigma/2-1}+x^{(1-\sigma)/2-1}$ times $e^{-\pi x}$ locally uniformly.
The right side of \eqref{eq:xi} is invariant under $s\mapsto1-s$, which is the functional equation for $\xi$.
Define $\zeta(s)=\pi^{s/2}\xi(s)/\Gamma(s/2)$; by Corollary~\ref{cor:gamma-nonzero} the factor $1/\Gamma(s/2)$ is entire, so $\zeta$ is meromorphic on $\C$, and it agrees with the Dirichlet series on $\re s>1$, hence with Lemma~\ref{lem:continuation} on $\re s>0$ by Lemma~\ref{lem:mero}(e).
\end{proof}

\begin{corollary}[Symmetry of the zeros]\label{cor:symmetry}
For $0<\re s<1$: $\zeta(s)=0\iff\zeta(1-s)=0$.
\end{corollary}
\begin{proof}
If $0<\re s<1$ then $\re(s/2)>0$ and $\re((1-s)/2)>0$, so $\Gamma(s/2)$ and $\Gamma((1-s)/2)$ are finite and nonzero (Corollary~\ref{cor:gamma-nonzero}), as are the powers of $\pi$.
Apply Theorem~\ref{thm:FE}.
\end{proof}

\section{A lower bound for \texorpdfstring{$\psi$}{psi}}\label{sec:cheb}

\begin{lemma}\label{lem:psi-lcm}
For every integer $n\ge1$, $\psi(n)=\log\lcm(1,2,\dots,n)$.
\end{lemma}
\begin{proof}
By Theorem~\ref{thm:uf}, $\lcm(1,\dots,n)=\prod_{p\le n}p^{k_p}$ where $k_p=\max_{j\le n}v_p(j)=\max\{k:p^{k}\le n\}$, and $\log$ of this is $\sum_{p\le n}k_p\log p=\sum_{p^{k}\le n}\log p=\psi(n)$.
\end{proof}

\begin{lemma}\label{lem:nair}
For integers $1\le m\le n$, $m\binom nm$ divides $\lcm(1,\dots,n)$. Consequently $\psi(2m)\ge 2m\log2-\log3$ for $m\ge1$, and $\psi(x)\ge x\log2-3$ for all real $x\ge2$.
\end{lemma}
\begin{proof}
Let $I=\int_0^1 v^{m-1}(1-v)^{n-m}\dd v$. Expanding $(1-v)^{n-m}$ by the binomial theorem,
$I=\sum_{j=0}^{n-m}\binom{n-m}{j}\frac{(-1)^{j}}{m+j}$, so $\lcm(1,\dots,n)\cdot I$ is an integer; it is positive since the integrand is positive.
By Lemma~\ref{lem:beta} with $a=m$, $b=n-m+1$ and Lemma~\ref{lem:gamma-basic}, $I=\frac{(m-1)!\,(n-m)!}{n!}=\frac{1}{m\binom nm}$.
Hence $\lcm(1,\dots,n)\cdot I=\lcm(1,\dots,n)/\bigl(m\binom nm\bigr)$ is a positive integer, which is the divisibility claim.
Take $n=2m$: $\lcm(1,\dots,2m)\ge m\binom{2m}{m}\ge m\cdot\frac{4^{m}}{2m+1}\ge\frac{4^{m}}{3}$,
because $4^{m}=\sum_{j=0}^{2m}\binom{2m}{j}\le(2m+1)\binom{2m}{m}$ and $m/(2m+1)\ge1/3$.
By Lemma~\ref{lem:psi-lcm}, $\psi(2m)\ge2m\log2-\log3$.
For real $x\ge2$ let $m=\lfloor x/2\rfloor\ge1$; then $\psi(x)\ge\psi(2m)\ge(x-2)\log2-\log3\ge x\log2-3$.
\end{proof}

\begin{lemma}\label{lem:Y-lower}
Let $Y(t)=tA(t)$ and $c_0=\frac{\log 2}{2e}$, $t_0=\frac{\log2}{6}$. Then $Y(t)\ge c_0$ for all $0<t\le t_0$.
\end{lemma}
\begin{proof}
For $0<t\le\tfrac12$, $A(t)\ge\sum_{n\le1/t}\Lam(n)e^{-nt}\ge e^{-1}\psi(1/t)\ge e^{-1}\bigl(\tfrac{\log2}{t}-3\bigr)$ by Lemma~\ref{lem:nair}, so $Y(t)\ge e^{-1}(\log2-3t)\ge e^{-1}\cdot\tfrac{\log2}{2}$ for $t\le t_0$.
\end{proof}

\section{Forward direction: \texorpdfstring{$\RH\Rightarrow\GR$}{RH implies G}}\label{sec:forward}

Throughout this section assume $\RH$. By Lemmas~\ref{lem:euler} and~\ref{lem:mertens}, $\RH$ is equivalent to
\begin{equation}\label{eq:RHprime}
\zeta(s)\ne0\quad\text{for all }s\text{ with }\re s>\tfrac12 .
\end{equation}

\begin{lemma}[Mellin representation of $A$]\label{lem:mellin-A}
For $t>0$ and $c>1$:
$\displaystyle A(t)=\frac1{2\pi}\int_\R\Gamma(c+i\tau)\Bigl(-\frac{\zeta'}{\zeta}\Bigr)(c+i\tau)\,t^{-c-i\tau}\dd\tau$.
\end{lemma}
\begin{proof}
By Corollary~\ref{lem:cahen-mellin} with $x=nt$, $e^{-nt}=\frac1{2\pi}\int_\R\Gamma(c+i\tau)n^{-c-i\tau}t^{-c-i\tau}\dd\tau$.
Multiply by $\Lam(n)$ and sum. The interchange of sum and integral is justified by Fubini, since
$\sum_n\Lam(n)n^{-c}\int_\R|\Gamma(c+i\tau)|\dd\tau<\infty$ by Lemmas~\ref{lem:euler} and~\ref{lem:gamma-decay}.
The inner sum is $-\zeta'/\zeta(c+i\tau)$ by Lemma~\ref{lem:euler}.
\end{proof}

\begin{lemma}[Borel--Carath\'eodory]\label{lem:BC}
Let $f$ be holomorphic on an open set containing the closed disc $|z|\le R$, with $\re f(z)\le M$ for $|z|\le R$.
Then for $0<r<R$ and $|z|\le r$,
\[
|f'(z)|\le\frac{4(R+r)}{(R-r)^{2}}\bigl(M-\re f(0)\bigr).
\]
\end{lemma}
\begin{proof}
Let $g=f-f(0)$ and $M_0=M-\re f(0)\ge\re g(0)=0$.
If $M_0=0$ then $\re g\le0=\re g(0)$ on the disc $|z|<R$, so $g$ is constant there by Lemma~\ref{lem:maxmod}(a), and the bound is trivial.
If $M_0>0$, put $\phi=g/(2M_0-g)$; the denominator has real part $\ge M_0>0$, and $|g|\le|2M_0-g|$ because this is equivalent to $\re g\le M_0$. So $\phi$ is holomorphic near the closed disc, maps it into the closed unit disc, and $\phi(0)=0$; Lemma~\ref{lem:maxmod}(b) gives $|\phi(z)|\le|z|/R$.
Then $|g(z)|\le\frac{|z|}{R}(2M_0+|g(z)|)$, so $|g(z)|\le\frac{2M_0|z|}{R-|z|}$.
On $|w|\le R_1=\frac{R+r}{2}$ this gives $|g(w)|\le\frac{2M_0(R+r)}{R-r}$, and Cauchy's estimate (Lemma~\ref{lem:cauchy-disc}) on the disc of radius $\frac{R-r}{2}$ about $z$ ($|z|\le r$) gives $|g'(z)|\le\frac{2M_0(R+r)}{R-r}\cdot\frac{2}{R-r}$.
\end{proof}

\begin{lemma}[Conditional bound for $\zeta'/\zeta$]\label{lem:cond-bound}
Assume \eqref{eq:RHprime} and fix $\varepsilon\in(0,\tfrac12)$. Then for $\tfrac12+\varepsilon\le\sigma\le2$ and $|\tau|\ge4$,
\[
\Bigl|\frac{\zeta'}{\zeta}(\sigma+i\tau)\Bigr|\le\frac{144}{\varepsilon^{2}}\log(|\tau|+2).
\]
\end{lemma}
\begin{proof}
Fix $\tau$ with $|\tau|\ge4$ and let $R'=\tfrac32-\tfrac\varepsilon4$, $R=\tfrac32-\tfrac\varepsilon2$, $r=\tfrac32-\varepsilon$.
Every $s$ with $|s-(2+i\tau)|<R'$ has $\re s>\tfrac12+\tfrac\varepsilon4$ and $|\im s|>|\tau|-\tfrac32\ge\tfrac52$, so by \eqref{eq:RHprime} and Lemma~\ref{lem:continuation} $\zeta$ is holomorphic and nonvanishing on this open disc.
By Lemma~\ref{lem:log} there is a holomorphic $f$ on $|z|<R'$ with $e^{f(z)}=\zeta(2+i\tau+z)$ and $f(0)=L(2+i\tau)$ (Lemma~\ref{lem:euler}), so $|f(0)|\le1$.
For $|z|\le R$ the point $s=2+i\tau+z$ has $\tfrac12\le\re s\le4$ and $2\le|\im s|\le|\tau|+2$, so by \eqref{eq:growth}
$\re f(z)=\log|\zeta(s)|\le\log12+\tfrac12\log(|\tau|+2)=:M$.
Lemma~\ref{lem:BC} gives, for $|z|\le r$,
\[
|f'(z)|\le\frac{4(R+r)}{(R-r)^{2}}\bigl(M-\re f(0)\bigr)\le\frac{4(R+r)}{(R-r)^{2}}(M+1).
\]
Here $R-r=\varepsilon/2$ and $R+r\le3$, so $4(R+r)/(R-r)^{2}\le48/\varepsilon^{2}$.
Since $|\tau|\ge4$, $\log(|\tau|+2)\ge\log6>1.79$, hence
$M+1=\log12+1+\tfrac12\log(|\tau|+2)\le\bigl(\tfrac{\log 12+1}{1.79}+\tfrac12\bigr)\log(|\tau|+2)\le3\log(|\tau|+2)$.
Therefore $|f'(z)|\le\frac{144}{\varepsilon^{2}}\log(|\tau|+2)$ for $|z|\le r$.
For $\tfrac12+\varepsilon\le\sigma\le2$ the point $z=\sigma-2$ satisfies $|z|\le\tfrac32-\varepsilon=r$, and $f'(z)=\zeta'/\zeta(\sigma+i\tau)$.
\end{proof}

\begin{theorem}\label{thm:forward}
Assume $\RH$. For every $\varepsilon\in(0,\tfrac12)$ there is $K_\varepsilon$ such that
$|A(t)-1/t|\le K_\varepsilon t^{-1/2-\varepsilon}$ for all $t>0$.
Consequently $\GR$ holds: $|r(t)|\le C_\varepsilon t^{1/2-\varepsilon}$ for $0<t\le\tfrac12$.
\end{theorem}
\begin{proof}
Fix $\varepsilon$ and put $c_1=\tfrac12+\varepsilon$, $c=2$, $h(s)=\Gamma(s)\bigl(-\zeta'/\zeta\bigr)(s)t^{-s}$.
On the open set $\re s>\tfrac12$ the function $h$ is meromorphic (Lemma~\ref{lem:mero}) with a single pole, at $s=1$ (Lemma~\ref{lem:continuation}, \eqref{eq:RHprime}, Corollary~\ref{cor:gamma-nonzero}), simple, with residue $\Gamma(1)\cdot1\cdot t^{-1}=1/t$.
Apply Lemma~\ref{lem:residue} to the rectangle with vertices $c_1\pm iT$, $c\pm iT$ for $T\ge4$:
\[
\int_{c-iT}^{c+iT}h\dd s-\int_{c_1-iT}^{c_1+iT}h\dd s+\int_{c_1}^{c}\bigl(h(\sigma-iT)-h(\sigma+iT)\bigr)\dd\sigma=\frac{2\pi i}{t}.
\]
On the horizontal edges, Lemmas~\ref{lem:gamma-decay} and~\ref{lem:cond-bound} give
$|h(\sigma\pm iT)|\le\frac{K}{1+T^{2}}\cdot\frac{144}{\varepsilon^{2}}\log(T+2)\cdot\max(t^{-c},t^{-c_1})\to0$ as $T\to\infty$.
By Lemma~\ref{lem:mellin-A} the first integral tends to $2\pi i\,A(t)$.
On the line $\re s=c_1$, $|\zeta'/\zeta(c_1+i\tau)|$ is bounded by $\frac{144}{\varepsilon^2}\log(|\tau|+2)$ for $|\tau|\ge4$ and by its maximum on the compact segment $|\tau|\le4$ (where it is continuous: no zero and no pole lies on $\re s=c_1$), so with Lemma~\ref{lem:gamma-decay}
$\int_\R|h(c_1+i\tau)|\dd\tau\le 2\pi K_\varepsilon\,t^{-c_1}$ for a constant $K_\varepsilon$ independent of $t$.
Letting $T\to\infty$: $A(t)=\frac1t+E(t)$ with $E(t)=\frac1{2\pi}\int_\R h(c_1+i\tau)\dd\tau$ and $|E(t)|\le K_\varepsilon t^{-1/2-\varepsilon}$.

For $0<t\le1$ this gives $|Y(t)-1|\le K_\varepsilon t^{1/2-\varepsilon}$ with $Y=tA$, hence by Lemma~\ref{lem:square}
$|\G(t)-1|=|Y(t)-1|\,|Y(t)+1|\le K_\varepsilon(2+K_\varepsilon)t^{1/2-\varepsilon}$, and for $0<t\le\tfrac12$,
$|r(t)|\le|\G(2t)-1|+|\G(t)-1|\le (1+\sqrt2)K_\varepsilon(2+K_\varepsilon)\,t^{1/2-\varepsilon}$.
\end{proof}

\section{Reverse direction: \texorpdfstring{$\GR\Rightarrow\RH$}{G implies RH}}\label{sec:reverse}

Throughout this section assume $\GR$. Fix $\varepsilon\in(0,\tfrac12)$, write $\alpha=\tfrac12-\varepsilon$, and let $C_\varepsilon,t_\varepsilon$ be as in $\GR$.

\subsection*{Step 1: the dyadic limit of $Y=tA$}

Let $c_0,t_0$ be as in Lemma~\ref{lem:Y-lower} and $t_1=\min(t_0/2,\,t_\varepsilon)$.
For $0<t\le t_1$ both $Y(t)$ and $Y(2t)$ are at least $c_0$, so by Lemma~\ref{lem:square}
\[
|Y(2t)-Y(t)|=\frac{|\G(2t)-\G(t)|}{Y(2t)+Y(t)}\le\frac{C_\varepsilon}{2c_0}\,t^{\alpha}.
\]
Hence for $0<t\le t_1$ and $K\ge1$,
\[
\bigl|Y(t)-Y(t2^{-K})\bigr|\le\sum_{k=1}^{K}\bigl|Y(t2^{-k+1})-Y(t2^{-k})\bigr|\le\frac{C_\varepsilon}{2c_0}t^{\alpha}\sum_{k\ge1}2^{-k\alpha}=C'\,t^{\alpha},
\qquad C'=\frac{C_\varepsilon}{2c_0(2^{\alpha}-1)} .
\]
The sequence $K\mapsto Y(t2^{-K})$ is therefore Cauchy; let $\ell(t)$ be its limit. Then for $0<t\le t_1$:
\begin{enumerate}
\item[(i)] $|Y(t)-\ell(t)|\le C't^{\alpha}$;
\item[(ii)] $\ell(t/2)=\ell(t)$, since both are the limit of the same sequence;
\item[(iii)] $\ell(t)\ge c_0$, as a limit of values $\ge c_0$;
\item[(iv)] $\ell$ is continuous on $(0,t_1]$: $Y$ is continuous (the series for $A$ converges uniformly on compact subsets of $(0,\infty)$), and $Y(t2^{-K})\to\ell(t)$ uniformly on $(0,t_1]$ because the tail of the telescoping sum is at most $C'(t2^{-K})^{\alpha}\le C't_1^{\alpha}2^{-K\alpha}$;
\item[(v)] $\ell$ is bounded on $(0,t_1]$, by (ii) and continuity on $[t_1/2,t_1]$.
\end{enumerate}
Consequently, for $0<t\le t_1$,
\begin{equation}\label{eq:A-ell}
\Bigl|A(t)-\frac{\ell(t)}{t}\Bigr|=\frac{|Y(t)-\ell(t)|}{t}\le C'\,t^{\alpha-1}.
\end{equation}

\subsection*{Step 2: meromorphic continuation of $\Gamma\cdot(-\zeta'/\zeta)$ to $\re s>\tfrac12+\varepsilon$}

For $\re s>1$, by Tonelli and $\int_0^\infty e^{-nt}t^{s-1}\dd t=\Gamma(s)n^{-s}$,
\begin{equation}\label{eq:mellin}
F(s):=\Gamma(s)\Bigl(-\frac{\zeta'}{\zeta}\Bigr)(s)=\int_0^\infty A(t)\,t^{s-1}\dd t .
\end{equation}
Split the integral at $t_1$ and insert $\ell$:
\[
F(s)=\underbrace{\int_0^{t_1}\Bigl(A(t)-\frac{\ell(t)}{t}\Bigr)t^{s-1}\dd t}_{I_1(s)}
+\underbrace{\int_0^{t_1}\ell(t)\,t^{s-2}\dd t}_{I_2(s)}
+\underbrace{\int_{t_1}^\infty A(t)\,t^{s-1}\dd t}_{I_3(s)}\qquad(\re s>1),
\]
each piece converging absolutely for $\re s>1$ (by \eqref{eq:A-ell}, by (v), and by $A(t)\le\sum_n ne^{-nt}=e^{-t}(1-e^{-t})^{-2}$).

$I_1$ is holomorphic on $\re s>1-\alpha=\tfrac12+\varepsilon$: by \eqref{eq:A-ell} the integrand is bounded by $C't^{\alpha+\sigma-2}$, integrable on $(0,t_1]$ when $\sigma>1-\alpha$, locally uniformly in $\sigma$ (Corollary~\ref{cor:morera}(c)).

$I_3$ is entire: for $t\ge t_1$, $A(t)\le e^{-t}(1-e^{-t_1})^{-2}$, so the integrand is dominated locally uniformly by $e^{-t}t^{\sigma-1}$ (Corollary~\ref{cor:morera}(c)).

$I_2$ is computed exactly from the dyadic invariance (ii). For $\re s>1$, substituting $t=2^{-k}u$ on $[t_12^{-k-1},t_12^{-k}]$,
\[
I_2(s)=\sum_{k\ge0}\int_{t_12^{-k-1}}^{t_12^{-k}}\ell(t)t^{s-2}\dd t
=\sum_{k\ge0}2^{-k(s-1)}\,\Phi(s)
=\frac{\Phi(s)}{1-2^{1-s}},\qquad
\Phi(s)=\int_{t_1/2}^{t_1}\ell(u)u^{s-2}\dd u ,
\]
where $\Phi$ is entire (continuous integrand on a compact interval, Corollary~\ref{cor:morera}(c)) and the geometric series converges since $|2^{1-s}|<1$.
The right side is meromorphic on $\C$ (Lemma~\ref{lem:mero}(b)), with poles contained in the set
\[
Z=\{\,1+2\pi ik/\log2:\ k\in\Z\,\}=\{s:2^{1-s}=1\},
\]
all of which lie on the line $\re s=1$.

Therefore $I_1+\Phi/(1-2^{1-s})+I_3$ is meromorphic on $\Omega_\varepsilon=\{\re s>\tfrac12+\varepsilon\}$ with poles in $Z$, and it agrees with $F$ on $\re s>1$.
Since $F=\Gamma\cdot(-\zeta'/\zeta)$ is itself meromorphic on $\re s>0$ (Lemmas~\ref{lem:gamma-basic} and~\ref{lem:continuation}), Lemma~\ref{lem:mero}(e) gives:
\begin{equation}\label{eq:poles}
\text{every pole of }F\text{ in }\Omega_\varepsilon\text{ lies in }Z,\ \text{hence on }\re s=1 .
\end{equation}

\subsection*{Step 3: no zeros of $\zeta$ in $\re s>\tfrac12$}

Suppose $\zeta(\rho)=0$ with $\rho\in\Omega_\varepsilon$, of multiplicity $m\ge1$. Then $\rho\ne1$, and by Lemma~\ref{lem:continuation} $-\zeta'/\zeta$ has a pole at $\rho$ with residue $-m\ne0$.
By Corollary~\ref{cor:gamma-nonzero}, $\Gamma(\rho)\ne0$ and $\Gamma$ is holomorphic at $\rho$ ($\re\rho>0$), so $F=\Gamma\cdot(-\zeta'/\zeta)$ has a pole at $\rho$ (Lemma~\ref{lem:mero}(d)).
By \eqref{eq:poles}, $\re\rho=1$. This contradicts Lemma~\ref{lem:mertens}.
Hence $\zeta$ has no zero in $\Omega_\varepsilon$. Since $\varepsilon\in(0,\tfrac12)$ was arbitrary, $\zeta$ has no zero with $\re s>\tfrac12$.

\subsection*{Step 4: conclusion}

Let $\rho$ be a zero of $\zeta$ with $0<\re\rho<1$. By Step 3, $\re\rho\le\tfrac12$. By Corollary~\ref{cor:symmetry}, $1-\rho$ is also a zero with $0<\re(1-\rho)<1$, so $\re(1-\rho)\le\tfrac12$, i.e.\ $\re\rho\ge\tfrac12$. Hence $\re\rho=\tfrac12$. This is $\RH$. \qed

\begin{remark}[Finite data suffice at each scale]\label{rem:finite}
Since $0\le\Lam(n)\le\log n$, $R(N)\le N\log^{2}N$. Let $0<t\le e^{-3}$, $\Lambda_t=\log(1/t)\ge3$ and $L=\lceil t^{-1}\Lambda_t\rceil$.
The function $x\mapsto x\log^{2}x\,e^{-xt}$ is decreasing for $x\ge L$ (its logarithmic derivative is $1/x+2/(x\log x)-t<0$ once $tx\ge3$ and $x\ge e$), so
\[
t^{2}\sum_{N>L}R(N)e^{-Nt}\le t^{2}\!\int_{L}^\infty\! x\log^{2}x\,e^{-xt}\dd x
=\int_{Lt}^\infty\! y\log^{2}(y/t)\,e^{-y}\dd y
\le2\!\int_{\Lambda_t}^\infty\!\bigl(y^{2}+y\Lambda_t^{2}\bigr)e^{-y}\dd y
\le 14\,t\,\Lambda_t^{3},
\]
using $\log^{2}(y/t)\le2(\log^{2}y+\Lambda_t^{2})$, $\log^{2}y\le y$ for $y>0$, and
$\int_{\Lambda}^\infty(y^{2}+y\Lambda^{2})e^{-y}\dd y=(\Lambda^{3}+2\Lambda^{2}+2\Lambda+2)e^{-\Lambda}\le7\Lambda^{3}e^{-\Lambda}$ for $\Lambda\ge1$.
This is $o(t^{1/2-\varepsilon})$ for every $\varepsilon<\tfrac12$.
So $\GR$ is unchanged if $\G(t)$ is replaced by $t^{2}\sum_{N\le L}R(N)e^{-Nt}$: at each scale $t$ the proposition is decided, up to an error strictly below the critical rate, by the finite list $R(4),\dots,R(L)$.
\end{remark}

\section{Dependency map}\label{sec:map}

\begin{center}
\begin{tabular}{p{0.34\textwidth}p{0.6\textwidth}}
Section \ref{sec:base} & the base: (B0) numbers and $\exp$, (B1) Riemann integration and uniform convergence, (B2) Lebesgue convergence theorems and Fubini \\
Section \ref{sec:complex} & Goursat, primitives, Cauchy's formula and estimates, Morera, Weierstrass, parameter integrals, identity theorem, maximum principle, Schwarz, logarithms, meromorphic functions, rectangles, residues \\
Section \ref{sec:arith} & B\'ezout, Euclid, unique factorisation, $\lcm$ \\
Section \ref{sec:gamma} & $\Gamma$: functional equation, Beta integral, reflection, nonvanishing, vertical decay \\
Section \ref{sec:fourier} & Gaussian transform, Fourier inversion, Fourier series, Poisson summation, Cahen--Mellin \\
Section \ref{sec:zeta} & Euler product, continuation to $\re s>0$, growth, $3$--$4$--$1$, theta inversion, functional equation, zero symmetry \\
Section \ref{sec:cheb} & $\psi(n)=\log\lcm(1..n)$, Nair's bound, $tA(t)\ge c_0$ \\
Section \ref{sec:forward} & Mellin inversion, Borel--Carath\'eodory, contour shift: $\RH\Rightarrow\GR$ \\
Section \ref{sec:reverse} & dyadic limit, exact dyadic Mellin sum, pole bookkeeping: $\GR\Rightarrow\RH$
\end{tabular}
\end{center}

Neither direction uses the prime number theorem, zero-density estimates, the Hadamard product, or any bound on $R(N)$ beyond the trivial $R(N)\le N\log^{2}N$.
The reverse direction does not need the limit $\ell$ to be the constant $1$; only its existence, positivity and dyadic invariance are used, and these follow from $\GR$ together with the elementary lower bound of Lemma~\ref{lem:nair}.

\end{document}
